\chapter{जनन-तंत्र}\label{ch:g8:reproductive-systems}

\Cref{ch:g8:sexual-reproduction} को \omterm{def:g8:sexual-reproduction:gametes}{युग्मक} चाहिए थे;
\cref{ch:g8:puberty} ने उनके कारख़ाने चालू कर दिए। यह अध्याय ख़ुद उन
कारख़ानों का नक़्शा बनाता है: \omterm{def:g4:animal-reproduction:sexual}{नर} और \omterm{def:g4:animal-reproduction:sexual}{मादा} जनन-तंत्र --- उनके अंग, हर एक
के बनाए \omterm{def:g8:sexual-reproduction:gametes}{युग्मक}, और \omterm{def:g4:animal-reproduction:sexual}{मादा} तंत्र की वह उल्लेखनीय मासिक लय। सादी शरीर-रचना,
सादे नामों के साथ: ये अंग भी \omterm{def:g4:heart-and-blood:heart}{दिल} और फेफड़ों जैसे ही अंग हैं, और यह उनका
अध्याय है।

\section{नर तंत्र}

\begin{definition}[नर जनन-तंत्र]\label{def:g8:reproductive-systems:male}
\omterm{def:g4:animal-reproduction:sexual}{नर} तंत्र के अंग:
\begin{itemize}
  \item दो \emph{वृषण}\index{वृषण}, जो शरीर की गरमाहट से बाहर
        \emph{वृषणकोश} में रहते हैं --- यानी \omterm{def:g8:sexual-reproduction:gametes}{युग्मकों} के कारख़ाने:
        \omterm{def:g8:puberty:puberty}{यौवनारंभ} के बाद से वे लगातार, करोड़ों की तादाद में
        \emph{शुक्राणु}\index{शुक्राणु} बनाते हैं;
  \item वे नलियाँ, जो शुक्राणुओं को जमा करती और ले जाती हैं, और रास्ते
        में वे ग्रंथियाँ भी जुड़ती हैं जिनके तरल उन्हें खिलाते तथा ढोते
        हैं --- और ये सब मिलकर \emph{वीर्य} बनाते हैं;
  \item \emph{शिश्न}\index{शिश्न}, जिससे होकर वीर्य शरीर से बाहर जाता
        है --- और वही अंग अपनी दूसरी ड्यूटी पर \omterm{def:g7:removing-wastes:kidneys}{मूत्र} भी ले जाता है
        (\cref{def:g7:removing-wastes:kidneys}); \omterm{def:g7:heart-and-circulation:rooms}{कपाटों} का एक इंतज़ाम
        दोनों ड्यूटियों को सख़्ती से अलग रखता है।
\end{itemize}
\end{definition}

\begin{example}[शुक्राणु, पास से]\label{ex:g8:reproductive-systems:sperm}
\omterm{def:g5:first-look-at-cells:microscope}{सूक्ष्मदर्शी} (\cref{def:g5:first-look-at-cells:microscope}) के नीचे
\omterm{def:g8:reproductive-systems:male}{शुक्राणु} सफ़र के लिए हलकी की हुई एक \omterm{def:g5:first-look-at-cells:cell}{कोशिका} है: एक \omterm{def:g1:my-body:parts}{सिर}, जो क़रीब-क़रीब
पूरा का पूरा \omterm{def:g6:cell-unit-of-life:parts}{केंद्रक} है --- यानी पिता का योगदान, कसकर बँधा हुआ
(\cref{rem:g5:first-look-at-cells:nucleus}) --- और एक फटकारती हुई
पूँछ। शरीर की सबसे छोटी \omterm{def:g5:first-look-at-cells:cell}{कोशिका}, तैरने के लिए बनी, और कुछ दिनों भर का
साज़-सामान लिए। उसका बनना \omterm{def:g8:puberty:puberty}{यौवनारंभ} से लेकर बुढ़ापे तक चलता है: \omterm{def:g4:animal-reproduction:sexual}{नर}
कारख़ाना लगातार चलने वाला है।
\end{example}

\section{मादा तंत्र}

\begin{definition}[मादा जनन-तंत्र]\label{def:g8:reproductive-systems:female}
\omterm{def:g4:animal-reproduction:sexual}{मादा} तंत्र के अंग, और सब पेट के निचले हिस्से में:
\begin{itemize}
  \item दो \emph{अंडाशय}\index{अंडाशय} --- यानी \omterm{def:g8:sexual-reproduction:gametes}{युग्मकों} के भंडार और
        कारख़ाने: \omterm{def:g8:reproductive-systems:male}{वृषणों} के उलट, वे आगे बनने वाले \emph{\omterm{def:g4:animal-reproduction:sexual}{अंडाणुओं}} का
        अपना पूरा भंडार जन्म से ही लिए रहते हैं, और \omterm{def:g8:puberty:puberty}{यौवनारंभ} के बाद
        उन्हें एक-एक करके छोड़ते हैं;
  \item दो \emph{अंडवाहिनियाँ}\index{अंडवाहिनी}, यानी झालरदार नलियाँ,
        जो छोड़े गए \omterm{def:g4:animal-reproduction:sexual}{अंडाणु} को लपक लेती हैं --- और यही
        \cref{prop:g8:sexual-reproduction:core} के दूसरे क़दम की आम
        मिलन-जगह है;
  \item \emph{\omterm{prop:g5:human-reproduction:plan}{गर्भाशय}} (\cref{prop:g5:human-reproduction:plan} वाला
        \omterm{prop:g5:human-reproduction:plan}{गर्भाशय}), जिसकी मोटी पेशीदार दीवार और हर बार नई होती परत
        \omterm{prop:g8:sexual-reproduction:core}{युग्मनज} का ठिकाना हैं;
  \item और \emph{योनि}, जो \omterm{prop:g5:human-reproduction:plan}{गर्भाशय} को बाहर से जोड़ती है --- जहाँ वीर्य
        लिया जाता है, और जो \cref{ex:g5:human-reproduction:birth} वाली
        जन्म की नली भी है।
\end{itemize}
\end{definition}

\begin{example}[अंडाणु, पास से]\label{ex:g8:reproductive-systems:egg}
\omterm{def:g4:animal-reproduction:sexual}{अंडाणु} भूमिका के सिवा हर बात में \omterm{def:g8:reproductive-systems:male}{शुक्राणु} का उलटा है: शरीर की
\emph{सबसे बड़ी} \omterm{def:g5:first-look-at-cells:cell}{कोशिका} --- एक मिलीमीटर का दसवाँ हिस्सा, यानी \omterm{def:g1:the-five-senses:sense}{आँख} की
पकड़ के किनारे पर --- गोल, बिना हिली-डुली, और सफ़र के पहले दिनों का
साज़-सामान लिए हुए (\cref{rem:g2:animals-grow-up:egg} वाला ज़र्दी का
उसूल, \omterm{def:g5:first-look-at-cells:cell}{कोशिका} के पैमाने पर)। एक कारख़ाना लगातार लाखों चलाता है; दूसरा
महीने में एक ही साज़-सामान वाली \omterm{def:g5:first-look-at-cells:cell}{कोशिका} छोड़ता है: यानी
\cref{prop:g4:animal-reproduction:strategies} वाला
संख्या-बनाम-देखभाल, ख़ुद \omterm{def:g8:sexual-reproduction:gametes}{युग्मकों} में लिखा हुआ।
\end{example}

\section{चक्र}

\begin{proposition}[मासिक चक्र]\label{prop:g8:reproductive-systems:cycle}
\omterm{def:g8:puberty:puberty}{यौवनारंभ} से लेकर लगभग पचास तक \omterm{def:g4:animal-reproduction:sexual}{मादा} तंत्र मोटे तौर पर \emph{28 दिन} का
एक दोहराता हुआ कार्यक्रम चलाता है --- यानी \emph{मासिक
चक्र}\index{मासिक चक्र}, जिसका हुक्म \cref{def:g8:puberty:hormones}
वाली हॉर्मोन-कड़ी देती है:
\begin{enumerate}
  \item \omterm{prop:g5:human-reproduction:plan}{गर्भाशय} की परत ख़ुद को दोबारा बनाती है, मोटी और \omterm{prop:g4:journey-of-food:absorb}{रक्त} से भरी
        --- यानी ठिकाना तैयार;
  \item चक्र के बीच में, चौदहवें दिन के आस-पास, कोई \omterm{def:g8:reproductive-systems:female}{अंडाशय} एक \omterm{def:g4:animal-reproduction:sexual}{अंडाणु}
        छोड़ता है --- यानी \emph{अंडोत्सर्ग}\index{अंडोत्सर्ग};
        \omterm{def:g8:reproductive-systems:female}{अंडवाहिनी} उसे लपक लेती है, और लगभग एक दिन तक वह निषेचित हो
        सकता है;
  \item \omterm{def:g5:flower-to-fruit:steps}{निषेचन} न हो: तब उस परत की ज़रूरत नहीं रहती; वह टूटकर कुछ दिनों
        में योनि से बाहर चली जाती है --- यही \emph{मासिक
        धर्म}\index{मासिक धर्म} है --- और कार्यक्रम फिर से शुरू हो
        जाता है;
  \item \omterm{def:g5:flower-to-fruit:steps}{निषेचन} हो जाए: तब \omterm{prop:g8:sexual-reproduction:core}{युग्मनज} उस तैयार परत में जम जाता है
        (\cref{def:g5:human-reproduction:pregnancy}), चक्र ठहर जाता
        है, और \omterm{def:g5:human-reproduction:pregnancy}{गर्भावस्था} शुरू हो जाती है।
\end{enumerate}
\end{proposition}
\admitted

\begin{omfigure}
\begin{tikzpicture}[scale=0.94]
  % cycle wheel
  \draw[very thick, omDef!60] (0,0) circle (2.3);
  % day marks
  \foreach \a/\d in {90/1, 0/8, -90/14, 180/22} {
    \fill[omDef] (\a:2.3) circle (0.07);
  }
  \node[font=\footnotesize, above] at (90:2.45) {दिन 1};
  \node[font=\footnotesize, right] at (0:2.45) {दिन 8};
  \node[font=\footnotesize, below] at (-90:2.45) {दिन 14};
  \node[font=\footnotesize, left] at (180:2.45) {दिन 22};
  % phases
  \draw[very thick, omThm] (90:2.3) arc[start angle=90, end angle=35,
    radius=2.3];
  \node[font=\footnotesize, omThm, align=center] at (58:3.5)
    {मासिक धर्म:\\ परत गिरती है};
  \node[font=\footnotesize, omMeth, align=center] at (-35:4.0)
    {परत दोबारा बनी,\\ मोटी और तैयार};
  \fill[omProp] (-90:2.3) circle (0.12);
  \node[font=\footnotesize, omProp, align=center] at (-90:3.55)
    {अंडोत्सर्ग:\\ एक अंडाणु छोड़ा गया};
  % center
  \node[font=\small, align=center] at (0,0.25)
    {चक्र:\\ लगभग 28 दिन};
  \node[font=\footnotesize, align=center] at (0,-0.75)
    {(निषेचन उसे ठहरा देता है:\\ गर्भावस्था)};
  % arrow direction
  \draw[->, thick, omDef] (150:2.3) arc[start angle=150,
    end angle=120, radius=2.3];
\end{tikzpicture}

{\small \omterm{prop:g8:reproductive-systems:cycle}{मासिक चक्र} एक पहिए की तरह: \omterm{prop:g8:reproductive-systems:cycle}{मासिक धर्म} उसे खोलता है, परत दोबारा
बनती है, \omterm{prop:g8:reproductive-systems:cycle}{अंडोत्सर्ग} बीच चक्कर में उसका ताज बनता है --- और \omterm{def:g5:flower-to-fruit:steps}{निषेचन} वह
अकेली घटना है जो पहिया रोक देती है।}
\end{omfigure}

\begin{remark}[चक्र के साथ जीना]\label{rem:g8:reproductive-systems:living}
व्यावहारिक बातें, सीधे-सीधे: पहले कुछ सालों तक चक्र अनियमित रहते हैं ---
मशीनरी सध रही होती है; लंबाई व्यक्ति-दर-व्यक्ति और महीने-दर-महीने बदलती
है; \omterm{prop:g8:reproductive-systems:cycle}{मासिक धर्म} में ऐंठन हो सकती है (\omterm{prop:g5:human-reproduction:plan}{गर्भाशय} एक \omterm{def:g3:bones-and-muscles:muscle}{पेशी} है,
\cref{def:g3:bones-and-muscles:muscle}, और गिरने के दौरान वह काम कर रही
होती है) --- और गरमाहट, हरकत तथा ज़रूरत पड़ने पर स्कूल की नर्स की सलाह,
सब मदद करते हैं; और पैड, टैम्पन तथा कप बस \omterm{prop:g8:reproductive-systems:cycle}{मासिक धर्म} का व्यावहारिक
साज़-सामान हैं, यानी हर दवाख़ाने की एक आम क़तार। इनमें से कुछ भी बीमारी
नहीं है: चक्र तो यह तंत्र है, जो ठीक अपनी बनावट के मुताबिक़ चल रहा है।
\end{remark}

\begin{method}[दोनों तंत्रों को आमने-सामने पढ़ना]\label{met:g8:reproductive-systems:sides}
नक़्शा साफ़ रखने के लिए काम के हिसाब से तुलना करो:
\begin{enumerate}
  \item \omterm{def:g8:sexual-reproduction:gametes}{युग्मक} का कारख़ाना: \omterm{def:g8:reproductive-systems:male}{वृषण} --- \omterm{def:g8:reproductive-systems:female}{अंडाशय};
  \item \omterm{def:g8:sexual-reproduction:gametes}{युग्मक}: लगातार तैरते लाखों --- या महीने में एक साज़-सामान
        वाला \omterm{def:g4:animal-reproduction:sexual}{अंडाणु};
  \item मिलने की जगह: \omterm{def:g8:reproductive-systems:female}{अंडवाहिनी}, जहाँ \omterm{def:g8:reproductive-systems:male}{शुक्राणु} योनि से \omterm{prop:g5:human-reproduction:plan}{गर्भाशय} होते
        हुए तैरकर पहुँचते हैं;
  \item ठिकाना: \omterm{def:g4:animal-reproduction:sexual}{नर} नक़्शे में कोई नहीं --- और \omterm{def:g4:animal-reproduction:sexual}{मादा} में हर महीने तैयार
        होता \omterm{prop:g5:human-reproduction:plan}{गर्भाशय};
  \item समय-सारणी: \omterm{def:g8:puberty:puberty}{यौवनारंभ} से लगातार --- या चक्रीय, \omterm{def:g8:puberty:puberty}{यौवनारंभ} से लगभग
        पचास तक।
\end{enumerate}
\end{method}

\section{अभ्यास}

\begin{exercise}[$\star$]\label{exo:g8:reproductive-systems:1}
\omterm{def:g4:animal-reproduction:sexual}{नर} तंत्र का नक़्शा बनाओ: हर अंग का नाम और उसका काम बताओ।
\end{exercise}

\begin{exercise}[$\star$]\label{exo:g8:reproductive-systems:2}
\omterm{def:g4:animal-reproduction:sexual}{मादा} तंत्र का नक़्शा भी वैसे ही बनाओ।
\end{exercise}

\begin{exercise}[$\star$]\label{exo:g8:reproductive-systems:3}
\omterm{def:g8:reproductive-systems:male}{शुक्राणु} और \omterm{def:g4:animal-reproduction:sexual}{अंडाणु} का फ़र्क़ बताओ: आकार, बनावट, संख्या, साज़-सामान।
\end{exercise}

\begin{exercise}[$\star$]\label{exo:g8:reproductive-systems:4}
\omterm{prop:g8:reproductive-systems:cycle}{अंडोत्सर्ग} पर क्या होता है, मोटे तौर पर किस दिन, और \omterm{def:g4:animal-reproduction:sexual}{अंडाणु} कितनी देर
तक निषेचित हो सकता है?
\end{exercise}

\begin{exercise}[$\star$]\label{exo:g8:reproductive-systems:5}
परत की भाषा में \omterm{prop:g8:reproductive-systems:cycle}{मासिक धर्म} क्या है --- और उसका आना उस चक्र के बारे में
क्या कहता है?
\end{exercise}

\begin{exercise}[$\star$]\label{exo:g8:reproductive-systems:6}
कौन-सी घटना चक्र को ठहरा देती है, और उसके बाद \omterm{prop:g5:human-reproduction:plan}{गर्भाशय} में क्या होता है?
\end{exercise}

\begin{exercise}[$\star\star$]\label{exo:g8:reproductive-systems:7}
\cref{prop:g8:sexual-reproduction:core} के तीनों क़दम इस अध्याय के
नक़्शे पर कहाँ-कहाँ होते हैं --- अंग-दर-अंग बताओ?
\end{exercise}

\begin{exercise}[$\star\star$]\label{exo:g8:reproductive-systems:8}
\omterm{def:g8:sexual-reproduction:gametes}{युग्मकों} को \cref{prop:g4:animal-reproduction:strategies} वाली दोनों
रणनीतियों के छोटे रूप की तरह पढ़ो। कौन-सा पक्ष कौन-सा तर्क ढोता है?
\end{exercise}

\begin{exercise}[$\star\star$]\label{exo:g8:reproductive-systems:9}
चक्रों के पहले कुछ साल अक्सर अनियमित क्यों रहते हैं, और लंबाई का बदलते
रहना सामान्य क्यों है? (\cref{rem:g8:reproductive-systems:living} में
से दो तथ्य।)
\end{exercise}

\begin{exercise}[$\star\star$]\label{exo:g8:reproductive-systems:10}
\omterm{prop:g8:reproductive-systems:cycle}{मासिक धर्म} की ऐंठन \cref{def:g3:bones-and-muscles:muscle} से समझाओ ---
और मदद करने वाली दो चीज़ों के नाम भी बताओ।
\end{exercise}

\begin{exercise}[$\star\star$]\label{exo:g8:reproductive-systems:11}
दोनों कारख़ाने उलटी समय-सारणियाँ चलाते हैं --- लगातार, बनाम चक्रीय और
जन्म से भंडार लिए। हर एक को उसके \omterm{def:g8:sexual-reproduction:gametes}{युग्मक} की बनावट तथा भूमिका से जोड़ो।
\end{exercise}

\begin{exercise}[$\star\star\star$]\label{exo:g8:reproductive-systems:12}
``चक्र वह सवाल है जो \omterm{prop:g5:human-reproduction:plan}{गर्भाशय} हर महीने पूछता है, और जिसका जवाब सिर्फ़
\omterm{def:g5:flower-to-fruit:steps}{निषेचन} दे सकता है।'' इस तस्वीर को
\cref{prop:g8:reproductive-systems:cycle} की चारों गिनी हुई घटनाओं में
खोलो।
\end{exercise}

\section{समस्या: दवाख़ाने के सिखाने वाले नमूने}

\begin{problem}[{सप्ताहांत समस्या --- दो नमूनों और एक कैलेंडर से तंत्र
समझाना}]\label{pb:g8:reproductive-systems:1}
किसी स्वास्थ्य केंद्र का खुला दिन तीन सिखाने वाली चीज़ें इस्तेमाल करता
है: \omterm{def:g4:animal-reproduction:sexual}{नर} तंत्र का एक नमूना, \omterm{def:g4:animal-reproduction:sexual}{मादा} तंत्र का एक नमूना, और 28 दिन के चक्र का
एक बड़ा कैलेंडर। समझाने वाले स्वयंसेवक तुम हो।

\medskip\noindent\textbf{भाग I --- दोनों नमूने।}
\begin{enumerate}
  \item एक आगंतुक \omterm{def:g4:animal-reproduction:sexual}{नर} नमूने पर \omterm{def:g8:reproductive-systems:male}{वृषणों} से लेकर बाहर निकलने तक की लकीर
        खींचता है। वह रास्ता सुनाओ, और साथ में बताओ कि रास्ते में क्या
        जुड़ता है तथा दोहरी ड्यूटी वाले अंग का इंतज़ाम क्या है।
  \item \omterm{def:g4:animal-reproduction:sexual}{मादा} नमूने पर दिखाओ कि \omterm{def:g8:sexual-reproduction:gametes}{युग्मकों} का भंडार कहाँ इंतज़ार करता है,
        \omterm{def:g5:flower-to-fruit:steps}{निषेचन} आम तौर पर कहाँ होता है, और \omterm{def:g5:human-reproduction:pregnancy}{गर्भावस्था} कहाँ ठहरती है।
  \item एक आगंतुक पूछता है कि नमूनों के युग्मक-कारख़ाने इतने तीखे रूप
        से अलग क्यों हैं --- रोज़ के लाखों बनाम महीने का एक।
        \cref{ex:g8:reproductive-systems:egg} के आख़िरी वाक्य से जवाब
        दो।
  \item कोई और पूछता है कि \omterm{def:g8:puberty:puberty}{यौवनारंभ} से पहले ये तंत्र कर क्या रहे थे।
        \cref{ch:g8:puberty} की मशीनरी से एक पंक्ति में जवाब दो।
\end{enumerate}

\medskip\noindent\textbf{भाग II --- कैलेंडर।}
\begin{enumerate}[resume]
  \item कैलेंडर पर चलो: \omterm{prop:g8:reproductive-systems:cycle}{मासिक धर्म}, दोबारा बनावट, \omterm{prop:g8:reproductive-systems:cycle}{अंडोत्सर्ग} और उपजाऊ
        दिन को उनकी मोटी-मोटी तारीख़ों के साथ रखो।
  \item एक आगंतुक मान बैठता है कि 28 दिन सबके लिए एक ही होते हैं। इस
        मान्यता को \cref{rem:g8:reproductive-systems:living} से
        सुधारो --- दो तरह की बदलाहटें बताकर।
  \item कैलेंडर पर दिखाओ कि चौदहवें दिन के आस-पास \omterm{def:g5:flower-to-fruit:steps}{निषेचन} हो जाए तो
        क्या बदलता है --- कौन-सी घटनाएँ रद्द होती हैं, कौन-सी शुरू
        (\cref{def:g5:human-reproduction:pregnancy} कहानी आगे बढ़ाती
        है)।
  \item एक किशोरी धीरे से पूछती है कि क्या ऐंठन का मतलब है कि कुछ
        ग़लत है। उसे \omterm{def:g3:bones-and-muscles:muscle}{पेशी} वाला जवाब और दोनों मदद बताओ --- और उनसे आगे
        की हर बात का पता भी।
\end{enumerate}

\medskip\noindent\textbf{भाग III --- समझाने वाले की कारीगरी।}
\begin{enumerate}[resume]
  \item आगंतुक नमूनों पर हँसने लगते हैं। तुम्हारा पहला वाक्य इन
        तंत्रों के साथ वैसा ही सुलूक करता है जैसा
        \cref{ch:g8:reproductive-systems} का पहला अनुच्छेद करता है।
        वह वाक्य लिखो।
  \item दोनों तंत्रों का सार
        \cref{met:g8:reproductive-systems:sides} वाली पाँच पंक्तियों
        की तुलना में, याददाश्त से लिखो।
  \item एक आगंतुक पूछता है कि शिशुओं की शुरुआत इन चीज़ों पर कहाँ बैठती
        है। तीनों चीज़ों को
        \cref{prop:g8:sexual-reproduction:core} के तीनों क़दमों में
        पिरो दो।
  \item खुले दिन का समापन एक वाक्य से करो: स्वास्थ्य केंद्र यह
        शरीर-रचना सिखाते ही क्यों हैं --- यानी गप के ख़िलाफ़ जानकारी,
        जैसा \cref{ex:g8:puberty:questions} ने कहा था।
\end{enumerate}
\end{problem}
